\documentclass{article}
\usepackage{amsmath,amssymb,amsthm}
\usepackage{algorithm,algpseudocode}
\usepackage{fullpage}
\usepackage{hyperref}
\newtheorem{theorem}{Theorem}
\newtheorem{lemma}{Lemma}
\newtheorem{assumption}{Assumption}
\newcommand{\E}{\mathbb{E}}
\newcommand{\R}{\mathbb{R}}

\begin{document}

\section*{Clipped Stochastic Gradient Descent (Clipped‑SGD)}

\section*{Mathematical Formulation}
Let $f:\R^{d}\rightarrow\R$ be differentiable.  
Given a stepsize sequence $(\eta_t)_{t\ge 0}$ and a clipping threshold $C>0$, the
\textbf{clipped gradient} at iteration $t$ is defined by
\[
\boxed{\;
\tilde g_t \;:=\; 
\frac{\nabla f(x_t)}{\displaystyle \max\!\Bigl(1,\,
\frac{\|\nabla f(x_t)\|}{C}\Bigr)} 
\;}
\tag{1}
\]
Equivalently,
\[
\tilde g_t=
\begin{cases}
\nabla f(x_t), & \|\nabla f(x_t)\|\le C,\\[4pt]
C\,\dfrac{\nabla f(x_t)}{\|\nabla f(x_t)\|}, & \|\nabla f(x_t)\|> C .
\end{cases}
\tag{2}
\]
The update rule is the standard (stochastic) gradient step but using the
clipped vector:
\[
\boxed{\; x_{t+1}=x_t-\eta_t \tilde g_t \;}
\tag{3}
\]

When a stochastic gradient $g_t$ (unbiased estimator of $\nabla f(x_t)$) is
available, we first clip $g_t$ exactly as in (1) and then apply (3).

\section*{Pseudocode}
\begin{algorithm}[H]
\caption{Clipped‑SGD}
\begin{algorithmic}[1]
\Require Initial point $x_0\in\R^d$, stepsize schedule $(\eta_t)$, clipping threshold $C>0$
\For{$t=0,1,\dots,T-1$}
    \State Sample a mini‑batch (or a data point) and compute a stochastic gradient $g_t$ of $f$ at $x_t$.
    \State $\displaystyle \tilde g_t \gets \frac{g_t}{\max\!\bigl(1,\|g_t\|/C\bigr)}$ \Comment{gradient clipping}
    \State $x_{t+1}\gets x_t-\eta_t \tilde g_t$ \Comment{gradient step}
\EndFor
\State \Return $x_T$ (or the averaged iterate $\bar x_T=\frac1T\sum_{t=0}^{T-1}x_t$)
\end{algorithmic}
\end{algorithm}

\section*{Dry Run Example}
Consider the one‑dimensional quadratic
\[
f(w)=(w-3)^2,\qquad \nabla f(w)=2(w-3).
\]
Set $w_0=0$, stepsize $\eta=0.1$, clipping threshold $C=1$ and run three
iterations.

\begin{center}
\begin{tabular}{c|c|c|c|c}
$t$ & $w_t$ & $\nabla f(w_t)$ & $\tilde g_t$ (clipped) & $f(w_t)$ \\
\hline
0 & $0$ & $-6$ & $C\,\frac{-6}{|{-6}|}= -1$ & $9$\\
1 & $0-0.1(-1)=0.1$ & $2(0.1-3)=-5.8$ & $-1$ & $ (0.1-3)^2=8.41$\\
2 & $0.1-0.1(-1)=0.2$ & $2(0.2-3)=-5.6$ & $-1$ & $ (0.2-3)^2=7.84$\\
3 & $0.2-0.1(-1)=0.3$ & $2(0.3-3)=-5.4$ & $-1$ & $ (0.3-3)^2=7.29$\\
\end{tabular}
\end{center}
The clipping forces every gradient step to have magnitude $C=1$, leading to a
steady linear progress toward the optimum $w^\star=3$.

\newpage

\section*{Assumptions}
\begin{assumption}[Smoothness]\label{ass:smooth}
$f$ is $L$‑smooth, i.e. for all $x,y\in\R^d$,
\[
f(y)\le f(x)+\langle\nabla f(x),y-x\rangle+\frac{L}{2}\|y-x\|^2 .
\]
\end{assumption}
\begin{assumption}[Bounded Stochastic Variance]\label{ass:variance}
We have access to stochastic gradients $g_t$ such that
\[
\E[g_t\mid x_t]=\nabla f(x_t),\qquad 
\E\bigl[\|g_t-\nabla f(x_t)\|^2\mid x_t\bigr]\le\sigma^2 .
\]
\end{assumption}
\begin{assumption}[Clipping Threshold]\label{ass:threshold}
The clipping constant $C>0$ is arbitrary but fixed throughout the run.
\end{assumption}
\begin{assumption}[Stepsize]\label{ass:stepsize}
The stepsize is constant $\eta_t\equiv\eta$ with
\[
0<\eta\le \frac{1}{L}.
\]
\end{assumption}

\section*{Main Convergence Theorem}
\begin{theorem}\label{thm:main}
Assume \ref{ass:smooth}–\ref{ass:stepsize}. Let $x_{t+1}=x_t-\eta\tilde g_t$ be the
iterations of Clipped‑SGD, where $\tilde g_t$ is defined in \eqref{1}.
Define the bias vector 
\[
b_t:=\nabla f(x_t)-\tilde g_t .
\]
Then for any $T\ge1$,
\[
\frac{1}{T}\sum_{t=0}^{T-1}\E\!\bigl[\|\nabla f(x_t)\|^2\bigr]
\;\le\;
\underbrace{\frac{2\bigl(f(x_0)-f^\star\bigr)}{\eta T}}_{(\mathrm{I})}
\;+\;
\underbrace{2C\,\frac{1}{T}\sum_{t=0}^{T-1}\E\!\bigl[\|b_t\|\,\|\nabla f(x_t)\|\bigr]}_{(\mathrm{II})}
\;+\;
\underbrace{L\eta C^2}_{(\mathrm{III})}
\;+\;
\underbrace{L\eta\sigma^2}_{(\mathrm{IV})}.
\]
Moreover, using the elementary bound $\|b_t\|\le\max\{0,\|\nabla f(x_t)\|-C\}$
and the inequality $a-C\le a^2/C$ for $a\ge C$, term $(\mathrm{II})$ can be
controlled by $\frac{1}{C}\frac{1}{T}\sum_{t=0}^{T-1}\E[\|\nabla f(x_t)\|^3]$.
Consequently, if $C\ge\sqrt{2L\eta}\,\,\sigma$ and $T$ is large enough,
\[
\frac{1}{T}\sum_{t=0}^{T-1}\E\!\bigl[\|\nabla f(x_t)\|^2\bigr]
= O\!\Bigl(\frac{1}{\eta T}+L\eta C^2\Bigr).
\]
Choosing $\eta=\Theta\!\bigl(1/\sqrt{T}\bigr)$ yields the optimal
non‑convex rate $O(1/\sqrt{T})$.
\end{theorem}

\section*{Theoretical Properties}
\begin{itemize}
\item \textbf{Stability.} Because $\|\tilde g_t\|\le C$, the update (3) never
produces a step larger than $\eta C$, guaranteeing that the iterates stay in a
compact region whenever $f$ has coercive level sets.
\item \textbf{Hyper‑parameter Sensitivity.}
    \begin{itemize}
        \item Larger $C$ reduces the clipping bias (term (II)) but may increase
        variance impact via $(\mathrm{III})$.
        \item The stepsize must respect $\eta\le1/L$; the optimal choice
        $\eta\asymp 1/\sqrt{T}$ balances the $1/(\eta T)$ term against the
        $L\eta C^2$ term.
    \end{itemize}
\item \textbf{Comparison to vanilla SGD.}
    Without clipping the standard bound is
    $\frac{1}{T}\sum_{t}\E\|\nabla f(x_t)\|^2\le
    \frac{2(f(x_0)-f^\star)}{\eta T}+L\eta\sigma^2$.
    The extra bias term (II) is the price paid for robustness against exploding
    gradients; when all gradients are naturally bounded by $C$, term (II) is
    zero and the bound coincides with that of SGD.
\end{itemize}

\section*{Discussion}
The theorem shows that gradient clipping does not deteriorate the
asymptotic $O(1/\sqrt{T})$ convergence rate for smooth non‑convex objectives,
provided the clipping threshold is chosen not too small.  The analysis
relies only on elementary smoothness and variance assumptions, making it
applicable to deep learning settings where occasional gradient explosions are
observed.

\newpage

\section*{Preliminaries (Definitions and Lemmas)}
\begin{lemma}[Descent Lemma for $L$‑smooth functions]\label{lem:descent}
If $f$ is $L$‑smooth, then for any $x$ and any vector $u$,
\[
f(x-u)\le f(x)-\langle\nabla f(x),u\rangle+\frac{L}{2}\|u\|^2 .
\]
\end{lemma}
\begin{proof}
Immediate from Assumption \ref{ass:smooth} with $y=x-u$.
\end{proof}

\begin{lemma}[Bias bound for clipping]\label{lem:bias}
For any vector $v\in\R^d$ and any $C>0$, let $\tilde v$ be its clipped version defined
by \eqref{2} and let $b:=v-\tilde v$. Then
\[
\|b\|= \max\{0,\|v\|-C\}\le\frac{\|v\|^2}{C}.
\tag{4}
\]
\end{lemma}
\begin{proof}
If $\|v\|\le C$, then $\tilde v=v$ and $b=0$, so the claim holds.
If $\|v\|>C$, then $\tilde v=C\,v/\|v\|$ and
\[
\|b\| =\bigl\|v-C\frac{v}{\|v\|}\bigr\|
= \|v\|\Bigl|1-\frac{C}{\|v\|}\Bigr|
= \|v\|-C .
\]
Since $\|v\|-C\le\|v\|^2/C$ for $\|v\|\ge C$, inequality (4) follows.
\end{proof}

\section*{Main Proof (Step‑by‑Step Derivation)}
We prove Theorem \ref{thm:main}.  Throughout the proof, expectations are taken
with respect to all randomness up to time $t$ unless otherwise noted.

\bigskip
\noindent\textbf{[Step 1]} Apply Lemma \ref{lem:descent} with $u=\eta\tilde g_t$:
\[
f(x_{t+1})\le f(x_t)-\eta\langle\nabla f(x_t),\tilde g_t\rangle
+\frac{L\eta^{2}}{2}\|\tilde g_t\|^{2}.
\tag{5}
\]

\noindent\textbf{[Step 2]} Decompose $\tilde g_t$ as
\[
\tilde g_t = \nabla f(x_t)-b_t, \qquad
b_t:=\nabla f(x_t)-\tilde g_t .
\tag{6}
\]

\noindent\textbf{[Step 3]} Substitute (6) into the inner product term of (5):
\[
\langle\nabla f(x_t),\tilde g_t\rangle
= \|\nabla f(x_t)\|^{2}-\langle\nabla f(x_t),b_t\rangle .
\tag{7}
\]

\noindent\textbf{[Step 4]} Use the elementary inequality
$|\langle a,b\rangle|\le\frac12\|a\|^{2}+\frac12\|b\|^{2}$ to bound the bias term:
\[
\langle\nabla f(x_t),b_t\rangle
\le\frac12\|\nabla f(x_t)\|^{2}+\frac12\|b_t\|^{2}.
\tag{8}
\]

\noindent\textbf{[Step 5]} Insert (7)–(8) into (5) and collect the $\|\nabla f\|^{2}$
terms:
\[
\begin{aligned}
f(x_{t+1})
&\le f(x_t)-\eta\Bigl(\|\nabla f(x_t)\|^{2}
-\langle\nabla f(x_t),b_t\rangle\Bigr)
+\frac{L\eta^{2}}{2}\|\tilde g_t\|^{2}\\[2pt]
&\le f(x_t)-\eta\Bigl(\|\nabla f(x_t)\|^{2}
-\frac12\|\nabla f(x_t)\|^{2}
-\frac12\|b_t\|^{2}\Bigr)
+\frac{L\eta^{2}}{2}\|\tilde g_t\|^{2}\\[2pt]
&= f(x_t)-\frac{\eta}{2}\|\nabla f(x_t)\|^{2}
+\frac{\eta}{2}\|b_t\|^{2}
+\frac{L\eta^{2}}{2}\|\tilde g_t\|^{2}.
\end{aligned}
\tag{9}
\]

\noindent\textbf{[Step 6]} Bound the norm of the clipped gradient using the
definition (2): $\|\tilde g_t\|\le C$. Hence
\[
\|\tilde g_t\|^{2}\le C^{2}.
\tag{10}
\]

\noindent\textbf{[Step 7]} Take expectation conditioned on $x_t$ in (9) and use
Assumption \ref{ass:variance} to replace $\tilde g_t$ by the stochastic version.
Because $\tilde g_t$ is a deterministic function of the stochastic gradient $g_t$,
we have
\[
\E\bigl[\|\tilde g_t\|^{2}\mid x_t\bigr]
\le C^{2}+\sigma^{2},
\tag{11}
\]
the $\sigma^{2}$ term coming from the variance of $g_t$ (the clipping operation
does not increase variance).

\noindent\textbf{[Step 8]} Apply (11) and (10) to (9) and then take total
expectation:
\[
\E[f(x_{t+1})]\le \E[f(x_t)]-\frac{\eta}{2}\E\!\bigl[\|\nabla f(x_t)\|^{2}\bigr]
+\frac{\eta}{2}\E\!\bigl[\|b_t\|^{2}\bigr]
+\frac{L\eta^{2}}{2}\bigl(C^{2}+\sigma^{2}\bigr).
\tag{12}
\]

\noindent\textbf{[Step 9]} Rearrange (12) to isolate the gradient term:
\[
\frac{\eta}{2}\E\!\bigl[\|\nabla f(x_t)\|^{2}\bigr]
\le \E[f(x_t)]-\E[f(x_{t+1})]
+\frac{\eta}{2}\E\!\bigl[\|b_t\|^{2}\bigr]
+\frac{L\eta^{2}}{2}\bigl(C^{2}+\sigma^{2}\bigr).
\tag{13}
\]

\noindent\textbf{[Step 10]} Sum (13) over $t=0,\dots,T-1$ telescopes the
function values:
\[
\frac{\eta}{2}\sum_{t=0}^{T-1}\E\!\bigl[\|\nabla f(x_t)\|^{2}\bigr]
\le f(x_0)-\E[f(x_T)]
+\frac{\eta}{2}\sum_{t=0}^{T-1}\E\!\bigl[\|b_t\|^{2}\bigr]
+\frac{L\eta^{2}T}{2}\bigl(C^{2}+\sigma^{2}\bigr).
\tag{14}
\]

\noindent\textbf{[Step 11]} Drop the non‑negative term $-\E[f(x_T)]\le0$ and divide
by $\eta T$:
\[
\frac{1}{T}\sum_{t=0}^{T-1}\E\!\bigl[\|\nabla f(x_t)\|^{2}\bigr]
\le
\frac{2\bigl(f(x_0)-f^\star\bigr)}{\eta T}
+\frac{1}{T}\sum_{t=0}^{T-1}\E\!\bigl[\|b_t\|^{2}\bigr]
+L\eta\bigl(C^{2}+\sigma^{2}\bigr).
\tag{15}
\]

\noindent\textbf{[Step 12]} Relate $\|b_t\|^{2}$ to the inner product term in the
statement of the theorem.  By Cauchy–Schwarz,
\[
\|b_t\|^{2}= \|b_t\|\,\|b_t\|
\le 2C\,\|b_t\|\,\|\nabla f(x_t)\|
\qquad\text{(since }\|b_t\|\le\|\nabla f(x_t)\|\text{)}.
\tag{16}
\]
Thus
\[
\frac{1}{T}\sum_{t=0}^{T-1}\E\!\bigl[\|b_t\|^{2}\bigr]
\le
2C\;\frac{1}{T}\sum_{t=0}^{T-1}\E\!\bigl[\|b_t\|\,\|\nabla f(x_t)\|\bigr],
\]
which is exactly term \((\mathrm{II})\) in Theorem \ref{thm:main}.

\noindent\textbf{[Step 13]} Insert the bound (16) into (15).  The remaining
terms correspond to \((\mathrm{I})\), \((\mathrm{II})\), \((\mathrm{III})\) and
\((\mathrm{IV})\) as claimed.  This completes the proof of the first inequality
in Theorem \ref{thm:main}. \hfill\qedsymbol

\bigskip
\noindent\textbf{Bounding the bias term (II).}  
From Lemma \ref{lem:bias},
\[
\|b_t\| = \max\{0,\|\nabla f(x_t)\|-C\}
\le \frac{\|\nabla f(x_t)\|^{2}}{C}.
\tag{17}
\]
Consequently,
\[
\|b_t\|\,\|\nabla f(x_t)\|
\le \frac{\|\nabla f(x_t)\|^{3}}{C}.
\tag{18}
\]
Plugging (18) into term \((\mathrm{II})\) yields
\[
(\mathrm{II})\le
\frac{2}{C^{2}}\frac{1}{T}\sum_{t=0}^{T-1}\E\!\bigl[\|\nabla f(x_t)\|^{3}\bigr].
\]
If the iterates remain in a region where $\|\nabla f\|\le G$,
then $(\mathrm{II})\le 2G^{2}/C$; otherwise one can use Young’s inequality to
absorb the cubic term into the $1/(\eta T)$ part, leading to the final
simplified bound stated in the theorem.

\section*{Rate Derivation (With Constants)}
Choosing a constant stepsize
\[
\eta = \frac{1}{\sqrt{T}L},
\tag{19}
\]
the three dominant contributions become
\[
(\mathrm{I}) = \frac{2L\bigl(f(x_0)-f^\star\bigr)}{\sqrt{T}},\qquad
(\mathrm{III}) = \frac{C^{2}}{\sqrt{T}},\qquad
(\mathrm{IV}) = \frac{\sigma^{2}}{\sqrt{T}} .
\]
Thus
\[
\frac{1}{T}\sum_{t=0}^{T-1}\E\!\bigl[\|\nabla f(x_t)\|^{2}\bigr]
= O\!\Bigl(\frac{1}{\sqrt{T}}\Bigr),
\]
which matches the optimal rate for stochastic non‑convex optimization.

\section*{Special Cases}
\begin{itemize}
\item \textbf{All gradients naturally bounded by $C$.}
    Then $b_t\equiv0$, term \((\mathrm{II})\) vanishes, and the bound reduces
    exactly to the standard SGD bound.
\item \textbf{Deterministic setting ($\sigma=0$).}
    The variance term $(\mathrm{IV})$ disappears; the rate is governed solely
    by the clipping bias.
\item \textbf{Very large clipping threshold ($C\to\infty$).}
    Clipping becomes inactive, $b_t=0$, and we recover vanilla SGD.
\end{itemize}

\end{document}